%%%%%%%%%%%%Outline%%%%%%%%%%%%%%%%%%

%%%
% message<Modern kernels increasingly execute deployment-specific programs on hot paths, so the real systems problem is no longer whether the kernel can run custom logic, but how that logic can be specialized and hardened after deployment.>
% deployment-specific networking, observability, security, and scheduling logic now lives on kernel fast paths
% once such logic is deployed, compilation becomes a lifecycle problem rather than a one-shot load-time step
% eBPF is the dominant concrete instance of this broader trend, not the whole story by itself

%%%
% message<The current compilation pipeline leaves opportunities at three different stages, and the common missing piece is a post-deployment compilation boundary.>
% LLVM to BPF misses deployment-specific and runtime-specific information
% BPF to native JIT lacks rich target-specific instruction selection and whole-program rewrite structure
% after deployment, the system cannot use profile feedback, stable runtime values, or new hardening opportunities to revisit already-loaded programs

%%%
% message<Neither a richer in-kernel optimizer nor load-time-only external optimization is the right answer.>
% pushing more policy into the verifier/JIT path puts fast-evolving complexity on a privileged and security-sensitive kernel boundary
% prior external optimizers improve initial code generation but cannot revisit live programs or exploit runtime information after deployment
% the missing abstraction is a principled split among kernel mechanism, userspace policy, and verifier-backed safety

%%%
% message<This paper's thesis is that programmable kernels need a post-deployment compilation boundary, not an in-kernel optimizing compiler.>
% principle 1: preserve the existing verifier safety boundary by re-verifying every rewritten program before replacement
% principle 2: place compilation policy in userspace so it can use richer analysis, runtime feedback, and rapid iteration
% principle 3: expose a minimal extensible low-level substrate for target-specific realization without growing a monolithic kernel optimizer

%%%
% message<\tool realizes this thesis by deliberately splitting responsibility so that each component stays narrow.>
% userspace daemon owns policy: profiling, rewrite construction, and deciding when specialization or hardening is worth attempting
% kernel recompilation component owns mechanism: recover original bytecode, route candidates through re-verification and re-JIT, and install them safely in place
% verifier owns safety and continues to reason about verifier-visible program semantics rather than trusting optimizer policy or native patches
% JIT and kinsn own low-level realization and map legal instruction variants to efficient native code on each platform

%%%
% message<This split is valuable because it improves generated code and adaptability without concentrating complexity inside the kernel.>
% the verifier does not learn profitability policy
% the JIT does not become a full runtime optimizer
% the kernel component does not grow into a general-purpose compiler
% userspace can evolve quickly without becoming part of the kernel's trusted safety boundary

%%%
% message<\tool turns eBPF from a one-shot verified load path into a post-deployment framework for specialization, hardening, and online adaptation.>
% it can revisit already-loaded programs through existing deployment paths
% it supports runtime-guided specialization as well as safety-oriented transformations
% replacement is in place, so applications, loaders, attachment paths, and steady-state execution remain unchanged

%%%
% message<Evaluation should promise a system evidence chain rather than only micro-optimization wins.>
% measure generated-code quality on microbenchmarks and real-program corpora
% measure end-to-end application benefit and online recompilation overhead
% measure safety preservation, fail-safe rejection, semantic correctness, and at least one safety/hardening case study
% measure extensibility of adding new capabilities

%%%
% message<Our contributions are a systems abstraction, a concrete realization, and a broad evidence set.>
% contribution 1: identify post-deployment kernel compilation as a systems problem and derive design principles
% contribution 2: design and implement \tool with responsibility split across userspace policy, kernel mechanism, verifier, and JIT/kinsn
% contribution 3: evaluate performance, code quality, online overhead, safety/correctness, and extensibility

%%%%%%%%%%%%Outline%%%%%%%%%%%%%%%%%%

\section{Introduction}
\label{sec:introduction}

Modern kernels increasingly execute deployment-specific programs on hot paths. These programs implement networking, observability, security, and scheduling policies that differ across machines and continue to evolve after deployment. Once such logic becomes part of the kernel fast path, compilation is no longer only a load-time concern; the system also needs a way to revisit, specialize, and harden deployed programs over time. Linux eBPF is the most prominent concrete instance of this trend. Production systems already rely on eBPF for programmable networking~\cite{cilium,katran}, observability~\cite{tracee}, security enforcement~\cite{krsi}, and scheduling~\cite{sched_ext}. The quality of compiled eBPF code therefore directly affects both system efficiency and the practicality of deployment-specific kernel services.

Today, however, the compilation boundary is in the wrong place. Opportunities are distributed across three stages of the pipeline. Before load time, LLVM-to-BPF compilation cannot see deployment-specific behavior or runtime-stable values. During BPF-to-native translation, the in-kernel JIT has only a narrow way to express target-specific instruction selection and whole-program rewrites. After deployment, the system still lacks a mechanism to use profile feedback, stable runtime values, or newly discovered hardening opportunities to revisit already-loaded programs. The problem is therefore not just that the kernel misses a few optimizations. The deeper problem is that already-deployed kernel programs have no post-deployment compilation boundary.

One response is to make the in-kernel verifier and JIT substantially smarter. Another is to optimize only before or during the initial load, as prior external eBPF optimization systems do~\cite{k2,merlin,epso}. Neither is sufficient. Building a rich optimizing compiler into the kernel would place fast-evolving, workload-sensitive, and platform-specific policy on one of the kernel's most privileged and security-sensitive paths. That choice raises implementation complexity, review cost, and the consequence of mistakes. Conversely, load-time systems can improve initial code generation, but they cannot revisit live programs or exploit the runtime information that appears only after deployment. What is missing is not another fixed pass pipeline, but a principled split among kernel mechanism, userspace policy, and verifier-backed safety.

We argue that programmable kernels need a post-deployment compilation boundary, not an in-kernel optimizing compiler. Such a boundary should satisfy three principles. First, it should preserve the existing verifier safety boundary: every rewritten program must be accepted through the normal verification path before replacement. Second, it should place compilation policy in userspace, where rewrite selection can use richer analysis, runtime feedback, and rapid iteration. Third, it should expose a minimal extensible low-level substrate through which rewritten programs can request target-specific realizations that plain BPF does not directly encode.

We present \tool, a post-deployment compilation framework for live eBPF programs that realizes these principles by deliberately splitting responsibility. A privileged userspace daemon owns policy: it profiles live programs, constructs whole-program rewrites, and decides when specialization or hardening is worth attempting. A small kernel recompilation component owns mechanism: it recovers the original bytecode of a live program, routes candidate rewrites back through verification and JIT compilation, and installs a new native image in place only after success. The verifier owns safety: it continues to reason about verifier-visible program semantics rather than trusting optimizer policy or native patches. The JIT and the \emph{kinsn} layer own low-level realization: they define how legal instruction variants map to efficient native code on each platform.

This split keeps each part narrow. The verifier does not learn profitability policy. The JIT does not become a full runtime optimizer. The kernel recompilation component does not grow into a general-purpose compiler. Meanwhile, userspace can evolve quickly without becoming part of the kernel's trusted safety boundary. The result is a system that can improve generated code and adapt to runtime behavior without concentrating complexity inside the kernel.

Together, these pieces turn eBPF from a one-shot verified load path into a post-deployment framework for specialization, hardening, and online adaptation. \tool can revisit already-loaded programs through existing deployment paths, apply runtime-guided specialization and target-specific rewrites, and support safety-oriented transformations such as hardening and fail-safe rejection of invalid rewrites. Because replacement occurs in place, existing applications, loaders, and attachment paths remain unchanged, and steady-state execution still follows the ordinary JIT fast path.

We evaluate \tool along five dimensions. First, we measure generated-code quality on controlled microbenchmarks and real-program corpora. Second, we measure end-to-end application benefit in realistic deployments. Third, we measure the online overhead of recompiling and replacing live programs. Fourth, we evaluate whether \tool preserves the existing eBPF safety model and the semantic correctness of rewrites, including fail-safe rejection and safety-oriented case studies. Finally, we evaluate extensibility by examining how new capabilities are added without turning the verifier or JIT into a large optimizing subsystem. Overall, the results show that post-deployment specialization can improve performance and code quality while keeping kernel-side complexity bounded.

\noindent\textbf{Contributions.}
In summary, we make three contributions:
\begin{itemize}
  \item We identify \emph{post-deployment kernel compilation} as a distinct systems problem and derive three design principles for it: kernel mechanism rather than kernel policy, verifier-backed safety rather than trusted rewrites, and an extensible low-level realization substrate rather than a monolithic in-kernel optimizer.
  \item We design and implement \tool, a post-deployment compilation framework that splits responsibility across userspace policy, a minimal kernel recompilation mechanism, the existing verifier, and an extensible JIT/\emph{kinsn} layer, enabling safe in-place recompilation of live eBPF programs.
  \item We evaluate \tool across microbenchmarks, real-program corpora, end-to-end deployments, online-overhead experiments, and safety/extensibility case studies, showing that it improves code quality and application performance while preserving the safety boundary and keeping kernel complexity small.
\end{itemize}
